% Generated from validated Article Draft Result. Do not hand-edit; revise the structured draft instead.
\section{Frontier Models — 6月、競争は「モデル性能」から提供形態へ広がった}
\noindent\textbf{2026年6月はMiniMax M3、Claude Sonnet 5、Qwen系更新、Gemini API、Kimi Codeが相次いで動いた。重要なのは一列の性能順位ではなく、open-weight、multimodal、API lifecycle、computer-use、coding/agent統合という提供形態の違いである。月末のGPT-5.6はlimited previewであり、7月の本格展開とは分けて読む必要がある。}\autocite{src-9191d43918fc,src-a24fd97eb3d0,src-d8e13386974c,src-eac8dfbeb78c,src-f6aa679babd7}

\subsection{6月の流れを、提供形態の変化として読む}

6月1日のMiniMax M3とAlibaba Model Studioの更新から始まり、Kimi CodeのGoal modeや/swarm、月後半のGemini APIのcomputer-use系更新、6月30日のAnthropicのモデル・科学・安全性関連発表へと、モデルそのものと利用面の更新が並行した。月末にはGPT-5.6 seriesのlimited previewも公表されたが、これは一般提供ではない。\autocite{src-30f48121decc,src-9191d43918fc,src-a24fd97eb3d0,src-d8e13386974c,src-eac8dfbeb78c,src-f6aa679babd7}


\subsection{open-weightとcomputer operationを同じモデル面に置くMiniMax M3}

MiniMaxはM3について、MSA sparse attention、最大1M context、画像・動画入力、desktop-computer operation、open weightsを組み合わせたモデルとして説明した。ここで注目すべきなのは、長文脈やmultimodalだけでなく、モデル配布とcomputer-use志向が一つの製品像にまとめられている点である。\autocite{src-d8e13386974c}


Anthropicの6月30日のnewsroomにはClaude Sonnet 5、Claude Science、Fable 5の再展開と安全性framework更新が同日に並ぶ。モデル更新を単独で見るより、agentic use、科学用途、安全性運用を一つのproduct lifecycleとして束ねていることが6月末の特徴だった。\autocite{src-f6aa679babd7}


Alibaba Model Studioの6月はQwen3.7-Plusの更新から始まり、Qwen3.7-Max snapshot、video系、image系へと複数modalのmodel lifecycleが続いた。さらにKimi K2.7 Codeの地域availabilityも記録され、model catalog自体が静的な一覧ではなく、能力・modal・地域提供を継続更新するsurfaceになっている。\autocite{src-eac8dfbeb78c}


Gemini APIでも、model shutdownやdeprecationと並行してstreaming speech、computer-use、月末の新preview/GAが進んだ。新モデルだけでなく、終了・移行・機能追加を含むAPI lifecycleが利用者の設計条件を変える。\autocite{src-a24fd97eb3d0}


Kimi Codeは6月にGoal mode、/swarm、Kimi K2.7 Codeのrelease/open-sourceを段階的に投入した。coding modelの更新を、単発のchat補助ではなくgoal-driven・multi-agentな開発workflowへ接続しようとする動きとして読める。\autocite{src-9191d43918fc}


Google DeepMind側でもDiffusionGemma、Gemma系更新、computer-useやvoice、agent-safety roadmapが6月の公式発信に並んだ。個別modelの優劣より、multimodal生成とagent操作、安全性設計が並行して開発面へ出てきたことが重要である。\autocite{src-dd5d9b8662f3}


\begin{claimboundary}
OpenAIが6月26日に示したGPT-5.6 seriesはlimited previewであり、Solは小規模なtrusted partner向けを含む限定段階だった。7月の一般提供を6月へ遡及させてはならない。6月号では、翌月の本格展開へつながるchronology bridgeとしてのみ位置付ける。\autocite{src-30f48121decc}
\end{claimboundary}

\subsection{同じ「frontier」でも、採用判断の物差しは異なる}

open-weightで持ち込めるモデル、managed APIとして更新され続けるモデル、coding/agent環境と一体化したモデルでは、利用者が評価すべきリスクと便益が違う。6月は「最高性能は何か」より、「どの実行環境へ、どの更新速度で、どの操作能力まで持ち込めるか」が競争軸として前景化した月だった。\autocite{src-9191d43918fc,src-a24fd97eb3d0,src-d8e13386974c,src-dd5d9b8662f3,src-eac8dfbeb78c}


\begin{claimboundary}
各社の公式release pageやnewsroomは、公開日・availability・記載機能を確認する一次資料である一方、capabilityやbenchmarkの優位性まで独立検証したものではない。本稿ではvendorの性能表現をそのまま横断ランキングへ変換しない。\autocite{src-9191d43918fc,src-a24fd97eb3d0,src-d8e13386974c,src-dd5d9b8662f3,src-eac8dfbeb78c,src-f6aa679babd7}
\end{claimboundary}

\begin{claimboundary}
GPT-5.6に関する6月時点のcapability、efficiency、performance表現もOpenAIのfirst-party claimとして扱う。limited previewというavailabilityの事実と、性能評価の主張は分離して読む。\autocite{src-30f48121decc}
\end{claimboundary}
